\chapter{Introduction}

This dissertation is concerned with the design of computing architectures which
support the implementation of \emph{declarative} programming languages. Programs
in declarative languages describe \emph{what} values are computed, unlike
programs in conventional imperative languages (such as PASCAL or ADA) which
describe \emph{how} values are computed. This fundamental shift in viewpoint is
made possible by rejecting the von-Neumann model of computation and adopting
models which are firmly rooted in mathematics.

Traditionally, there have been two distinct types of declarative language:
\emph{functional} and \emph{logic}. Functional languages are based on the
\lambdacalculus{} or recursive equation systems. Logic languages, on the other
hand, are based on a procedural interpretation of the first order predicate
calculus. In recent years it has become apparent that these two paradigms are
different aspects of a single underlying theory based on meaning-preserving
transformations, and the integration of functions and relations into a single
language is a very active area of research.

The mathematical foundations of declarative languages endow them with elegant
semantics and an expressiveness that is unsurpassed among programming languages.
Such power does not come for free, however, and declarative languages pay a stiff
price in terms of implementation efficiency on conventional von-Neumann
computers. In order to increase the efficiency of declarative languages,
implementors have gradually chipped away at their mathematical foundations,
eliminating features which are too expensive.

The research described here is an attempt to redress this situation, primarily
with regard to functional languages based on the \lambdacalculus. The work
described herein is motivated by the conviction that a full and complete
implementation of the \lambdabetacalculus{} is desirable from both a theoretical
and a practical standpoint. If conventional von-Neumann computers are not well
suited to implement the \lambdacalculus, why not design a computer that is? This
dissertation explores this idea, starting with the theory of the
\lambdacalculus{} and progressing to the design of a sequential computing
architecture which directly implements the \lambdabetacalculus{} as its native
machine language.\footnote{This work is based on an invention by Dr. Klaus J.
Berkling entitled ``Lambda Calculus Computer,'' U.S. Patent Application No.
247,622, filing date 22 September 1988. Inquiries regarding the
commercialization of the head order reduction technology presented herein should
be directed to Syracuse University.} Interestingly, the resulting architecture,
named \RED{2},\footnote{This architecture is the second in a series of reduction
machines that began with \RED{1}, developed by Berkling and his colleagues at
GMD~[32]. \RED{1} is a string reduction machine for the \lambdacalculus,
\RED{2} is a graph reduction machine for the \lambdacalculus, and \RED{3} is
planned to provide a complete integration of the \ensuremath{\lambda}- and
predicate calculi.} turns out to be quite close to conventional von-Neumann
technology.

The organization of this dissertation presents an ordered development of
\RED{2}, from theory to machine. Chapter 2 presents the theory of the
\lambdacalculus{} and some of the reasons why it is so difficult to implement.
An equivalent, operationally oriented theory is then presented as the basis for a
machine implementation. Chapter 3 introduces the syntax and operational semantics
of THOR, an extension of the \lambdacalculus{} which is the assembly language
for \RED{2}. The architecture of \RED{2} is presented in Chapter 4. A conscious
attempt is made in Chapters 2--4 to establish, through both formal and informal
arguments, that \RED{2} correctly implements the complete
\lambdabetacalculus. In Chapter 5 the performance of \RED{2} is evaluated with
respect to other implementation techniques and several enhancements to the
architecture are presented. Chapter 6, the final chapter, explores the
possibility of implementing the Horn clause subset of the predicate calculus on
\RED{2}.
